% Discussion

\section{Discussion}
\label{sec:discussion}

\subsection{Limitations}

\textbf{Detection Scope.}
\textsc{WasmSentinel} currently focuses on syscall-based detection, which may miss attacks that operate entirely within user space (e.g., some in-memory attacks).
Extending coverage to include memory access patterns, library calls, or hardware performance counters is future work.

\textbf{Evasion Attacks.}
Sophisticated adversaries may attempt to evade ML-based detection through:
\begin{itemize}
    \item \textit{Mimicry attacks}: Inserting benign syscalls to mask malicious patterns
    \item \textit{Polymorphic behavior}: Varying attack execution to avoid learned patterns
    \item \textit{Adversarial examples}: Crafted inputs that fool the neural network
\end{itemize}

While our hybrid architecture provides some resilience (STIDE and MLP may catch different evasion attempts), a determined adversary with knowledge of our detection algorithm could potentially evade detection.
Defense-in-depth with multiple detection layers remains important.

\textbf{Dataset Bias.}
Our evaluation uses the ADFA-LD dataset, which, while well-established, dates from 2013 and may not fully represent modern attack techniques.
Evaluation on more recent datasets (e.g., LID-DS, custom container workloads) is planned for future work.

\textbf{Model Generalization.}
The current model is trained on specific attack types.
Generalization to entirely new attack categories (zero-day attacks) may be limited.
Continuous model updates with new attack data are necessary for maintaining detection effectiveness.

\subsection{Security Analysis}

\textbf{Wasm Sandbox Escapes.}
The security of \textsc{WasmSentinel} relies on Wasmtime's sandbox guarantees.
Wasmtime has undergone extensive security review and implements:
\begin{itemize}
    \item Memory isolation through virtual memory mapping
    \item Control flow integrity via indirect call validation
    \item Resource limiting through fuel-based execution metering
\end{itemize}

To date, no sandbox escape vulnerabilities have been disclosed for Wasmtime in production use.

\textbf{eBPF Verifier Trust.}
We trust the Linux eBPF verifier to ensure safety of loaded eBPF programs.
The verifier has had historical bugs~\cite{ebpf-bugs}, but it is part of the kernel TCB for any eBPF-based security tool.

\textbf{Supply Chain Security.}
\textsc{WasmSentinel}'s architecture naturally supports supply chain security for detection rules:
\begin{itemize}
    \item Rules are content-addressed (can be signed and verified)
    \item Sandbox isolation limits damage from malicious rules
    \item Resource limits prevent denial-of-service
\end{itemize}

\subsection{Deployment Considerations}

\textbf{Kernel Version Requirements.}
\textsc{WasmSentinel} requires Linux 5.8+ for ring buffer support.
Older kernels can use perf buffers with reduced performance.

\textbf{Performance Tuning.}
For high-throughput environments, we recommend:
\begin{itemize}
    \item Increasing ring buffer size to prevent event loss
    \item Adjusting STIDE threshold for workload-specific false positive rates
    \item Pinning the detection thread to dedicated CPU cores
\end{itemize}

\textbf{Model Updates.}
The modular design allows hot-swapping of detection models without restarting the daemon.
We envision a workflow where:
\begin{enumerate}
    \item Security team trains updated models on new attack data
    \item New Wasm module is signed and distributed
    \item Running instances atomically swap to the new module
\end{enumerate}

\subsection{Future Work}

\textbf{Multi-Dataset Evaluation.}
We plan to evaluate on additional datasets including LID-DS, CICIDS2017, and custom container workload traces.

\textbf{Advanced ML Models.}
Exploring more sophisticated architectures:
\begin{itemize}
    \item Temporal convolutional networks (TCN) for sequence modeling
    \item Graph neural networks for process relationship analysis
    \item Transformer-based models for attention over syscall sequences
\end{itemize}

\textbf{Federated Learning.}
Enabling collaborative model updates across multiple deployments while preserving privacy of individual organizations' security telemetry.

\textbf{Hardware Acceleration.}
Investigating Wasm SIMD extensions and potential WebGPU integration for accelerated inference.
